% Ad Atticum 4.13 (ad-atticum-04-13) --- English translation
% Translated by Claude (Anthropic) from the Latin (Perseus).
% Style guide: STYLE.md.

\ciceroLetterOpener{CICERO ATTICO salutem}

\ciceroSection{1} I see you know we came to Tusculum on 17 K. Dec. There Dionysius was on hand for us. We mean to be at Rome on 14 K. Dec. — what do I say, mean? rather, are compelled. Milo's wedding. There is a certain expectation about the elections. I, on condition the elections actually hold, do not mind having been absent from the altercations which I hear took place in the senate. For I should either have defended a position that did not please me, or have failed someone I should not have failed. But by Hercules, I would have you write me as much as you can about those matters and the present state of the commonwealth and the spirit in which the consuls bear this \textit{[Greek: skylmon]} — the present harassment. I am sharp-set \textit{[Greek: oxypeinos]}; and, if you ask, everything is suspect to me.

\ciceroSection{2} Our Crassus, they say, has set out in his commander's cloak with less dignity than once L. Paulus did, his peer in being a consul a second time. O the worthless fellow! As for the books on oratory, the matter has been seen to carefully on my side. They have long been much in my hands. You may copy them. Once again I beg you: \textit{[Greek: tēn parousan]} — the lady of the house — that I do not come to you as a stranger.
